\scanpage{41}
1) Cas $n$ pair, $\rho=2\widetilde\rho$, $\sigma=\widetilde\sigma$, l'équation s'écrit

\[
\lambda\sigma(1+\deleted{\unsure{2}}\lambda\widetilde\rho)=0
\]

Si $\sigma$ \deleted{\unsure{est}}\inserted{et $\widetilde\rho$} inversibles, la seule solution \uline{inversible} (\inserted{i.e.} correspondant à une section de $U_\rho$ partout $\ne 1$) est donnée par

\[
\lambda=-\frac{1}{\deleted{\unsure{2|4}}\widetilde\rho}\qquad
\text{d'où}\qquad
\underbrace{u=\mathrm{id}-\frac{1}{\deleted{\unsure{2|4}}\widetilde\rho}\pi'\otimes\pi}_{\text{\uline{antipodisme}}}
\]

\deleted{Si \uline{$\sigma=0$} $\unsure{G_\rho|U_\rho|\mathcal U_\rho}$ invarie $f$, \unsure{ou|ce|on} ce qui revient au $=$,}

D'autre part, conditions équivalentes

\diagram{\begin{tikzcd}[ampersand replacement=\&,column sep=0pt,row sep=small]
{} \& {}\\
\mathrm b)\ \&U_\rho\ \text{invarie }f^E\\
\mathrm c)\ \&U_\rho\ \text{invarie } Q_X^{\deleted{\unsure{E|X}}}\\
\mathrm a)\arrow[uuu,bend left=75]\ \&\sigma=0
\end{tikzcd}}

Enfin, si $\widetilde\rho=0$, \deleted{deux}\inserted{ces} conditions \unsure{équivalent|équivalentes} encore à la suivante

\deleted{a) $\sigma=0$}

d) $\exists$ loc. section \deleted{\unsure{divise}} de $U_\rho$ partout $\ne1$ qui invarie $f^E$. \deleted{(ou, ce qui revient au $=$)}\inserted{(ou aussi, $Q_X$)}.
% archive sheet number: 40
